% Cote 119: offprints, reports, courses, CV and letters (typed and printed part): archive cote 119, site page 202 (PDF page 203).
% typescript. Draft compared: none.
% AI-transcribed by Claude Opus 5.5 (claude-opus-5-5, adaptive thinking, effort high), Message Batches API, batch r119_full (msgbatch_01DwZNm5Y7YsKBWyAN4sFs8w), 2026-11-13; not yet checked by eye.
\scanpage{119/202}
\mspage{- 7 -}

elle-même contenue dans la théorie des classes de Chern \hand{$\ell$}-adiques des fibrés
vectoriels sur des topos annelés quelconques. Dans [13 bis], j'établis, à peu
de choses près, que pour un groupe discret $G$ , la théorie des représentations
linéaires de $G$ (sur un anneau de base quelconque) ne dépend que du
complété profini'' $\hat{G}$ de $G$ .

8. \uline{Groupes formels} ([17] [16 SGA 7] [14 bis])

C'est un sujet qui relève à la fois de la théorie des groupes, de
celle des groupes de Lie, de la géométrie algébrique, de l'arithmétique, et
(sous la forme voisine des groupes de Barsotti-Tate) de la théorie des systèmes
locaux. Ici encore, la théorie des schémas permet une grande aisance, et c'est
dans ce contexte par exemple que se place d'emblée I. Manin [$*$a], dans son
exposé classique de la théorie de Dieudonné. Ma principale contribution, en
dehors de cette simplification conceptuelle, a été le développement d'une
``théorie de Dieudonné'' pour les groupes de Barsotti-Tate sur des schémas de
base généraux à caractéristiques résiduelles $>0$ , en termes du ``cristal de
Dieudonné'' associé à un tel groupe. Une esquisse de cette théorie a été expo-%
sée dans divers cours et séminaires, y compris dans mon cours au Collège de
France en 1970/71 et 71/72; certains énoncés principaux sont esquissés dans
les \unsure{C.R.|CR.} du Congrès International de Nice en 1970 [14 bis]. Une partie de ces
idées est développée dans la thèse de W. Messing [$*$], et les besoins techniques
de la théorie ont été la motivation pour le développement par L. Illusie [$*$]
de sa théorie des déformations des schémas en groupes commutatifs, vérifiant
des conjectures suggérées par cette ``théorie de Dieudonné cristalline''. Par
ailleurs, les relations entre schémas abéliens et groupes de Barsotti-Tate asso-%
ciés sont explorées et exploitées également dans [17] et dans [16,SGA 7,Exp.IX].

\uline{Références} : J. Tate, B. Mazur, A. Néron, L. Illusie, J.N. Katz, W. Messing,
I. Manin.

9. \uline{Arithmétique} ([16 SGA 5, Exp XVI] [18, Exp III])

Ma contribution principale a consisté (en collaboration avec M. Artin)
en la démonstration de la rationalité des fonctions L associées à des faisceaux
\hand{$\ell$}-adiques généraux sur des variétés algébriques sur des corps finis, compre-%
nant comme cas particulier les fonctions L associées à des caractères de groupes
finis opérant sur de telles variétés. S'inspirant des conjectures de Weil, on
arrive en effet à exprimer ces fonctions L en termes de produits alternés de

\margin{202}
